\begin{graphillustration}{Перестройка гамильтонова пути в цикл: первый участок проходим вперёд, второй — назад. Точками показано ребро пути, не вошедшее в цикл.}{fig:10-hamilton-rewire}
\begin{tikzpicture}[font=\small]
\node[gi small,label=below:{$u=x_1$}] (u) at (0,0) {};
\node[gi small] (a) at (1.1,0) {};
\node[gi small,label=below:{$x_i$}] (b) at (2.2,0) {};
\node[gi small,label=below:{$x_{i+1}$}] (c) at (3.3,0) {};
\node[gi small] (d) at (4.4,0) {};
\node[gi small,label=below:{$v=x_n$}] (v) at (5.5,0) {};
\draw[gi edge] (u) -- (a);
\draw[gi edge] (a) -- (b);
\draw[gi edge] (b) -- (c);
\draw[gi edge] (c) -- (d);
\draw[gi edge] (d) -- (v);
\draw[gi emphasis] (u) to[bend left=40] (c);
\draw[gi alternate] (b) to[bend right=40] (v);
\node[gi note] at (2.75,1.3) {Путь и два ребра при $i\in S\cap T$};
\begin{scope}[shift={(0,-3)}]
\node[gi small,label=below:{$u=x_1$}] (u) at (0,0) {};
\node[gi small] (a) at (1.1,0) {};
\node[gi small,label=below:{$x_i$}] (b) at (2.2,0) {};
\node[gi small,label=below:{$x_{i+1}$}] (c) at (3.3,0) {};
\node[gi small] (d) at (4.4,0) {};
\node[gi small,label=below:{$v=x_n$}] (v) at (5.5,0) {};
\draw[gi arrow] (u) -- (a);
\draw[gi arrow] (a) -- (b);
\draw[gi arrow] (v) -- (d);
\draw[gi arrow] (d) -- (c);
\draw[gi arrow] (b) to[bend right=40] (v);
\draw[gi arrow] (c) to[bend right=40] (u);
\draw[gi removed] (b) -- (c);
\end{scope}
\node[gi note] at (8,-1.5) {$u\leadsto x_i\to v$\\$v\leadsto x_{i+1}\to u$\\[5pt]Ребро $x_ix_{i+1}$ пропускаем.};
\end{tikzpicture}
\end{graphillustration}
